%% cap02-marco-conceptual.tex — Marco conceptual
%% Rejilla de comparación: los siete apartados de la sección 6 de Miao et al. (2021).
%% Sustituye a cap02-marco-teorico.tex (seis ejes confucianos) cuando el cambio de
%% marco quede aprobado. Mientras tanto no se incluye en main.tex.

\providecommand{\pendiente}[1]{\textbf{[Pendiente: #1]}}

\chapter{Marco conceptual}
\label{cap:marco-conceptual}

Este capítulo tiene dos partes que responden a preguntas distintas. El estado del arte
(sección~\ref{sec:estado-arte}) reúne las investigaciones recientes que compararon
políticas de inteligencia artificial y, en particular, las que se ocuparon de la educación:
qué problemas estudiaron, cómo los definieron, con qué evidencia y a qué resultados
llegaron. El marco teórico (sección~\ref{sec:marco-teorico}) expone las posturas desde las
cuales analizo esas políticas: la educación comparada, los marcos internacionales de
inteligencia artificial y educación, y la competencia docente en inteligencia artificial.
Sigo en esta división a \citet{londono2014guia}, que distinguen el estado del arte, dirigido
a dar cuenta de las investigaciones recientes sobre las categorías de análisis, del marco
teórico, que analiza las posturas disciplinares sobre esas mismas categorías.

%% ─────────────────────────────────────────────────────────────────────────────
\section{Antecedentes: estado del arte}
\label{sec:estado-arte}

\subsection{El análisis comparado de estrategias nacionales de inteligencia artificial}
\label{subsec:ea-estrategias}

La primera ola de estudios comparados no trató las estrategias nacionales como tales, sino
los documentos de principios éticos que gobiernos, empresas y organismos publicaron a partir
de 2016. \citet{jobin2019landscape} reunieron 84 documentos con principios o directrices
éticas para la inteligencia artificial y encontraron una convergencia en torno a cinco
principios: transparencia, justicia y equidad, no maleficencia, responsabilidad y
privacidad. La convergencia, sin embargo, era de vocabulario. Los autores documentaron una
divergencia sustantiva en cuatro frentes: cómo se interpreta cada principio, por qué se
considera importante, a qué asuntos y actores se aplica y cómo debería llevarse a la
práctica. Ningún principio aparecía en todo el corpus.

El mismo estudio dejó ver un sesgo de origen que importa para esta tesis. Estados Unidos y
Reino Unido aportaban más de un tercio de los documentos, África y América del Sur no
estaban representadas de forma independiente, y los autores advirtieron que los países
económicamente más desarrollados estaban dando forma a un debate que se presentaba como
global \citep{jobin2019landscape}. Su búsqueda se limitó a documentos en inglés, alemán,
francés, italiano y griego, y ellos mismos reconocieron que un sesgo de idioma pudo inclinar
el corpus hacia los resultados en inglés.

\citet{fjeld2020principled} llegaron a una conclusión parecida por otra vía. Con una
muestra intencional de 36 documentos influyentes, buscando la máxima variación entre
autores (gobiernos, organismos intergubernamentales, empresas, asociaciones profesionales
y organizaciones civiles), agruparon 47 principios en ocho temas: privacidad, rendición de
cuentas, seguridad, transparencia y explicabilidad, equidad y no discriminación, control
humano de la tecnología, responsabilidad profesional y promoción de valores humanos
\citep[pp.~4, 14--15]{fjeld2020principled}. Los documentos más recientes tendían a cubrir
los ocho temas, lo que sugería una convergencia; pero los autores advirtieron que cubrir más
temas no hace mejor a un documento, porque el contexto importa, y que entre los conceptos y
su puesta en práctica queda una brecha amplia \citep[pp.~5, 66]{fjeld2020principled}. Su
corpus incluyó documentos en inglés, chino, francés, alemán y español, y aun así no
encontraron ninguno de África \citep[pp.~14--15]{fjeld2020principled}. Dejaron fuera,
además, las estrategias nacionales que piden crear principios sin proponer ninguno
\citep[pp.~12--13]{fjeld2020principled}, de modo que su mapa de consensos éticos no
describe todavía lo que los gobiernos se proponen hacer.

La educación ocupaba un lugar marginal en estos estudios. En el corpus de
\citet{jobin2019landscape} no figura como principio; aparece como medio para generar
confianza en la tecnología o como propuesta de introducir la ética en la formación
científica y tecnológica. Cuando la guía de UNESCO para responsables de política revisó el
panorama, registró, con datos de la OCDE, más de trescientas iniciativas de política de inteligencia artificial en
sesenta países, la mayoría con alguna referencia a la educación, y las clasificó según
trataran la educación como política independiente, integrada en una estrategia general o
temática \citep[p.~28]{miao2021guidance}. El paso siguiente fue preguntar qué dicen
exactamente las estrategias nacionales cuando hablan de educación.

\subsection{La educación en las estrategias nacionales: el estudio de Schiff}
\label{subsec:ea-schiff}

\citet{schiff2022education} hizo esa pregunta de manera sistemática. Partió de 76
documentos del sector público, aplicó criterios de exclusión (documentos incompletos,
documentos de organismos intergubernamentales, documentos no disponibles en inglés y
duplicados de un mismo país) y llegó a una muestra de 24 estrategias nacionales publicadas
entre mediados de 2016 y principios de 2020, con cerca de 1{,}491 páginas en total
\citep[pp.~531--532, 552]{schiff2022education}. La muestra incluye a China y a México, este
último con el documento \textit{Towards an AI Strategy in Mexico}
\citep[p.~535]{schiff2022education}.

Su procedimiento es el antecedente metodológico más cercano a esta tesis. Construyó un
códebook preliminar de once categorías, lo probó sobre cinco documentos elegidos al azar y
lo depuró hasta siete códigos: cinco temas y dos categorías de contexto
\citep[pp.~533, 553--554]{schiff2022education}. Los cinco temas se agrupan en dos dominios.
El primero, \textit{educación para la inteligencia artificial}, reúne la formación de
expertos, la preparación de la fuerza laboral y la alfabetización pública en inteligencia
artificial. El segundo, \textit{inteligencia artificial para la educación}, reúne las
herramientas de enseñanza y aprendizaje y las herramientas administrativas
\citep[p.~534]{schiff2022education}. Cada documento se codificó por presencia o ausencia de
cada tema, con una categoría intermedia para los casos dudosos.

El resultado dio título al artículo. La educación era tema prioritario en 21 de las 24
estrategias, y las 24 trataban al menos uno de los temas de educación para la inteligencia
artificial; solo nueve abordaban con alguna profundidad la inteligencia artificial para la
educación, y solo cinco (China, India, Kenia, Malta y España) trataban los dos temas de ese
dominio \citep[p.~535]{schiff2022education}. Schiff concluyó que los responsables de
política ven la educación sobre todo como un instrumento para desarrollar la fuerza laboral
y formar expertos \citep[p.~550]{schiff2022education}. La ética de la inteligencia
artificial en educación aparecía tan poco que no alcanzó a constituirse como código
\citep[p.~534]{schiff2022education}.

El propio autor declaró tres límites que esta tesis atiende de forma directa. La
codificación la hizo un solo investigador, de modo que no fue posible estimar la
confiabilidad entre codificadores \citep[p.~533]{schiff2022education}. El corpus se
restringió a documentos en inglés por el dominio de idiomas del autor
\citep[p.~552]{schiff2022education}. Y las estrategias nacionales de inteligencia artificial
reflejan solo una parte del discurso sobre educación, porque los organismos educativos
pueden tener planteamientos más detallados que no aparecen en ellas
\citep[p.~550]{schiff2022education}. Su distinción entre educación para la inteligencia
artificial e inteligencia artificial para la educación organiza una de mis preguntas de
investigación.

\subsection{Rejillas de comparación aplicadas a políticas escolares de inteligencia artificial}
\label{subsec:ea-rejillas}

Si Schiff mostró qué lugar ocupa la educación en las estrategias generales de inteligencia
artificial, otros estudios compararon directamente las políticas educativas. \citet{kundu2025aipolicies}
analizaron cinco documentos de gobierno sobre la integración de la inteligencia artificial
en la educación escolar de China, Singapur, Finlandia y Estados Unidos. Para Estados Unidos
usaron dos documentos estatales, de California y Massachusetts, y no uno federal. Los
examinaron con los criterios SMART y con un marco de nueve componentes operativos que
construyeron los propios autores. Del análisis surgieron nueve proposiciones temáticas,
entre ellas la integración del currículo y la ética, la formación docente dinámica, el
financiamiento equitativo, las alianzas entre actores y el monitoreo adaptativo.
Encontraron que Finlandia y Singapur tienen políticas con un fuerte componente ético y
centrado en las personas, mientras que China y Estados Unidos se inclinan hacia la
innovación y el desarrollo de la fuerza laboral, y que la implementación enfrenta en los
cuatro casos brechas de equidad, falta de preparación docente y desajustes de
contexto.\footnote{Para este estudio solo tuve acceso al resumen publicado; la descripción
de sus nueve componentes deberá verificarse contra el texto completo.}

El trabajo demuestra que una rejilla de componentes sirve para comparar políticas
escolares de países con sistemas políticos muy distintos. Su límite, para los fines de esta
tesis, está en el origen de la rejilla: la construyeron los autores, de modo que otro
investigador con otros componentes podría llegar a otro ordenamiento de los países. Y su
caso estadounidense ilustra la advertencia de \citet{bray1995levels} sobre los sistemas
federales, porque lo que se compara de Estados Unidos no es una política nacional sino dos
estatales.

\subsection{El análisis comparado en español y en América Latina}
\label{subsec:ea-latam}

En la educación comparada de lengua española el tema es reciente y se ha concentrado en la
educación superior. \citet{martinezrojas2026regulacion} compararon las normas sobre
inteligencia artificial de seis universidades, tres de España (Universidad Autónoma de
Madrid, UNED y Universitat Oberta de Catalunya) y tres de Estados Unidos (Harvard, Stanford
y Pennsylvania), elegidas entre las más influyentes de cada país que ya contaban con
regulación específica \citep[pp.~423--424]{martinezrojas2026regulacion}. Partieron de
categorías previas tomadas de la Comisión Europea y de UNESCO, dos autores revisaron las
fuentes de manera independiente y el análisis produjo tres temas (recomendaciones
generales, de docencia y de investigación) con doce subtemas
\citep[pp.~425--426]{martinezrojas2026regulacion}. Encontraron un propósito común, reducir
los daños posibles sin impedir el uso, y una diferencia de énfasis: las universidades
españolas son más prescriptivas en la enseñanza y las estadounidenses dan más orientación
sobre investigación y datos \citep[pp.~431, 436]{martinezrojas2026regulacion}. Concluyeron
que la cultura institucional puede pesar más que la política nacional, y pidieron
políticas guiadas por necesidades pedagógicas y no solo por imperativos tecnológicos
\citep[pp.~435--436]{martinezrojas2026regulacion}. Advierten, además, que no buscaban
jerarquizar países ni universidades, sino entender cómo los factores locales moldean la
adopción de la inteligencia artificial \citep[p.~424]{martinezrojas2026regulacion}.

\citet{lara2026gobernanza} compararon México y Costa Rica con un enfoque cuantitativo
basado en datos oficiales secundarios de 2019 a 2023. Relacionaron el índice de desarrollo
humano con el gasto en educación superior, la matrícula terciaria, la proporción de
empresas que usan inteligencia artificial y un índice de cooperación entre universidad e
industria. Costa Rica presentó valores promedio superiores en todos los indicadores, y la
vinculación entre universidad y empresa fue la variable con mayor peso en su modelo. Su
tesis es que el principal desafío que enfrentan ambos países no es tecnológico sino
institucional, y que la inteligencia artificial amplifica las capacidades o las
limitaciones de los sistemas que ya existen. El estudio es valioso porque trata la
inteligencia artificial como objeto de gobernanza y porque incluye a México, pero su
modelo estadístico descansa en datos agregados de solo dos países durante cinco años, un
número de observaciones reducido para estimar cinco predictores, y sus propios autores
reconocen que la información agregada limita el análisis de las dinámicas institucionales.

Los dos trabajos comparten dos rasgos. Estudian dos casos cada uno, y se ocupan de la
educación superior y no de la política educativa nacional en su conjunto. Ninguno llega a
siete países ni emplea una rejilla tomada de un organismo internacional como esquema de
codificación.

\subsection{La crítica a la orientación de los organismos internacionales}
\label{subsec:ea-critica}

La rejilla que adopto no es neutral, y la literatura lo ha documentado.
\citet{mochizuki2025ethics} analizaron las dos guías de UNESCO sobre inteligencia
artificial en educación, la de responsables de política de 2021 y la de inteligencia
artificial generativa de 2023. Combinaron el análisis del discurso, siguiendo la pregunta
de Bacchi sobre cómo se representa el problema, con un análisis de redes de las 190 obras
que las dos guías citan y de sus autores \citep[pp.~2--4]{mochizuki2025ethics}.

Sus hallazgos se dirigen a la fuente misma de mi criterio de comparación. En la guía de
2021 la ética ocupa el primer plano, pero se despliega para presentar la regulación y la
formación de capacidades como condiciones previas para aprovechar los beneficios de la
tecnología, lo que evita preguntar si la inteligencia artificial es necesaria o deseable en
educación \citep[pp.~5, 8]{mochizuki2025ethics}. En el apartado 6.5, el problema que se
representa es un conjunto de deficiencias del sistema educativo sin relación con las
preocupaciones éticas del resto del documento, y el efecto es presentar como inevitable la
adopción de la inteligencia artificial y producir la figura de docentes y estudiantes que
necesitan desarrollar competencias \citep[p.~11]{mochizuki2025ethics}. El análisis de redes
mostró que la guía descansa sobre todo en literatura angloestadounidense, que entre las
afiliaciones más citadas figuran empresas como McKinsey, OpenAI y Google Health, y que el
consultor que participó en la redacción de ambas guías ocupa el centro de la red
\citep[pp.~12--15]{mochizuki2025ethics}. Los autores concluyen que, para UNESCO, la ética
de la inteligencia artificial y el tecno-solucionismo no compiten entre sí, y señalan la
ausencia de voces del Sur Global \citep[pp.~17--18]{mochizuki2025ethics}.

Esta crítica no invalida la rejilla como instrumento de comparación, pero obliga a declarar
de dónde viene. Uso los siete apartados como categorías para describir qué atiende y qué
omite cada política, no como un programa normativo que suscribo ni como el estándar que los
países deberían cumplir.

\subsection{Síntesis: coincidencias, diferencias, tendencias y contradicciones}
\label{subsec:ea-sintesis}

Leídos en conjunto, estos estudios coinciden en que las políticas de inteligencia
artificial pueden compararse con rejillas de categorías, y en que el vocabulario común
esconde diferencias de fondo. \citet{jobin2019landscape} y \citet{fjeld2020principled}
encontraron convergencia en los principios y divergencia en su interpretación y puesta en
práctica; \citet{schiff2022education} encontró que las estrategias hablan mucho de
educación, pero casi siempre como formación de fuerza laboral.

Difieren en el origen de la rejilla. Jobin y sus colaboradores, Fjeld y sus colaboradores y
Schiff derivaron sus categorías del propio corpus, por codificación inductiva o por
depuración de un códebook preliminar. \citet{kundu2025aipolicies} construyeron la suya
como autores. \citet{martinezrojas2026regulacion} partieron de categorías de la Comisión
Europea y de UNESCO, pero las aplicaron a normas universitarias y no a políticas
nacionales. Ninguno de los estudios revisados toma una lista cerrada de un organismo
internacional y la aplica como esquema de codificación a las políticas educativas de un
conjunto de países.

Se advierten tres tendencias. La primera es la comparación de pocos casos: dos en los
estudios de lengua española y cuatro en el de Kundu y Bej. La segunda es un sesgo de idioma
y de origen que los propios autores reconocen: corpus restringidos al inglés o a unas
cuantas lenguas europeas y ausencia de África y de buena parte de América Latina
\citep{jobin2019landscape, fjeld2020principled, schiff2022education}. La tercera es el
tránsito de los principios éticos generales hacia la política educativa concreta, primero
en la educación superior y después en la educación escolar.

La contradicción más visible está entre el uso creciente de los marcos de UNESCO como
referencia para diseñar y comparar políticas y la crítica documentada a su orientación
tecno-solucionista y a la base angloestadounidense y empresarial de sus fuentes
\citep{mochizuki2025ethics}. Esta tesis se ubica en ese punto: usa la rejilla de UNESCO
porque es la lista más completa y específica de educación disponible, y declara sus sesgos
en lugar de ignorarlos. Los vacíos que atiende son tres: comparar siete países en lugar de
dos a cuatro, incluir documentos en su lengua original y no solo en inglés, y medir el
acuerdo en la codificación, que Schiff no pudo estimar por haber trabajado como único
codificador.

%% ─────────────────────────────────────────────────────────────────────────────
\section{Marco teórico}
\label{sec:marco-teorico}

\subsection{La política educativa como objeto de comparación}
\label{subsec:mt-comparacion}

\subsubsection*{a) El método comparativo de Bereday}

El método comparativo clásico en educación organiza el trabajo en cuatro etapas:
descripción, interpretación, yuxtaposición y comparación. La comunidad de comparatistas lo
atribuye sobre todo a \citet{bereday1964comparative}, aunque la tradición alemana lo
atribuye a Franz Hilker, y ambos autores se citaron mutuamente en varias ocasiones
\citep{adick2018bereday}. Las cuatro etapas gozan de aceptación entre los comparatistas,
aunque desarrollos teóricos posteriores las matizan y amplían
\citep[p.~25]{senent2016iniciacion}.

La secuencia estructura esta tesis. En la descripción se presenta cada documento en su
propio contexto nacional. En la interpretación se leen los documentos a la luz de ese
contexto. En la yuxtaposición se colocan los casos lado a lado, en cuadros organizados por
criterios de comparación, como lo hacen \citet[p.~25]{senent2016iniciacion}; aquí esos
criterios son los siete apartados de la guía de UNESCO. En la comparación se formulan las
semejanzas y diferencias que la yuxtaposición hizo visibles. La rejilla de UNESCO entra,
por tanto, en la tercera etapa, y no sustituye a la lectura contextual de las dos primeras.

\subsubsection*{b) Los niveles de comparación de Bray y Thomas}

\citet{bray1995levels} criticaron que la educación comparada se hubiera concentrado en
países y regiones del mundo, porque esa concentración ignora las diferencias entre
estados, distritos, escuelas, aulas e individuos. Para ordenar el campo propusieron un
cubo de tres dimensiones. La primera es geográfica y tiene siete niveles: regiones del
mundo, países, estados o provincias, distritos, escuelas, aulas e individuos. La segunda
agrupa a la población por rasgos no geográficos, como etnia, religión, edad o género, o la
toma completa. La tercera recoge los aspectos de la educación y de la sociedad que se
comparan, entre ellos el currículo, los métodos de enseñanza, el financiamiento y las
estructuras de gestión. Todo estudio comparado, sostienen, ocupa una o varias celdas de ese
cubo \citep[p.~474]{bray1995levels}.

El marco me obliga a decir dónde se ubica este trabajo. Comparo en el nivel de país, sobre
la población completa y sin desagregar por grupo, y en el aspecto de la política de
inteligencia artificial en educación. Bray y Thomas piden a quien trabaja en un solo nivel
que declare los límites de esa elección \citep[p.~488]{bray1995levels}, y el límite aquí es
concreto: advierten que en Australia, Canadá, Alemania y Estados Unidos las autoridades
estatales o provinciales pueden controlar la estructura de la educación, al grado de que
resulta difícil hablar de un sistema educativo nacional \citep[p.~479]{bray1995levels}.
Cuatro de los siete países de esta tesis son federales en ese sentido. Por eso el documento
nacional de cada uno debe leerse como la declaración del gobierno federal, no como la
política que efectivamente rige en sus escuelas.

\subsubsection*{c) La política como texto, discurso y efecto}

\citet{rizvi2010globalizing} parten de la definición de Easton, para quien la política es
la asignación autoritativa de valores, y sostienen que las políticas
son normativas: expresan fines y medios para orientar la conducta de las personas
\citep[pp.~4, 7]{rizvi2010globalizing}. Retoman de Ball la idea de que la política es a la
vez texto y acción, lo que se enuncia y lo que se lleva a cabo, y de que la investigación
sobre políticas examina tres planos: los textos, los discursos que los enmarcan y sus
efectos \citep[pp.~4--5, 12]{rizvi2010globalizing}. Los textos de política, añaden, suelen
coser intereses en competencia, de modo que un documento puede reunir valores que no se
concilian del todo \citep[p.~6]{rizvi2010globalizing}.

Dos precisiones de estos autores delimitan lo que puedo afirmar. La primera es que una
política puede ser simbólica: se anuncia sin compromiso real de implementación y sin
financiamiento asignado \citep[p.~9]{rizvi2010globalizing}. La segunda es que los docentes
interpretan la política y la traducen a la práctica, de modo que lo escrito no llega
intacto al aula \citep[pp.~5--6]{rizvi2010globalizing}. Mi análisis opera en los planos del
texto y del discurso y no mide efectos. Lo que encuentre en un documento es lo que ese
gobierno declara, no lo que ocurre en sus escuelas.

Los mismos autores sitúan el análisis en un contexto de globalización: las políticas
nacionales se vinculan con discursos globalizados y con la presión de organismos como la
Unión Europea, la OCDE y la UNESCO, pero esas presiones se manifiestan de forma vernácula,
según la cultura, la historia y la política de cada país
\citep[pp.~x, xii]{rizvi2010globalizing}. Esa tensión entre un vocabulario común y su
traducción nacional es justamente lo que una comparación con rejilla internacional puede
hacer visible.

\subsection{Los marcos internacionales de inteligencia artificial y educación}
\label{subsec:mt-marcos}

\subsubsection*{a) El Consenso de Beijing}

El Consenso de Beijing es el documento final de la Conferencia Internacional sobre
Inteligencia Artificial y Educación que UNESCO organizó en mayo de 2019, con cincuenta
ministros y viceministros y cerca de quinientos representantes de más de cien Estados
miembros \citep[párr.~1]{unesco2019beijing}. Sus recomendaciones a los gobiernos se
organizan en diez áreas: la planeación de la inteligencia artificial en las políticas
educativas; la gestión y prestación de la educación; el empoderamiento de la enseñanza y de
los docentes; el aprendizaje y su evaluación; los valores y competencias para la vida y el
trabajo; el aprendizaje a lo largo de la vida; el uso equitativo e inclusivo; la equidad de
género; el uso ético, transparente y auditable de los datos y algoritmos educativos; y el
monitoreo, la evaluación y la investigación. Una undécima área, sobre financiamiento y
cooperación internacional, se dirige a los organismos internacionales
\citep[párrs.~8--38]{unesco2019beijing}.

El documento reafirma un enfoque humanista: la inteligencia artificial debe estar bajo
control humano y centrada en las personas \citep[párrs.~6--7]{unesco2019beijing}. Para esta
tesis importan dos de sus compromisos. Sobre los docentes, afirma que la interacción entre
docentes y estudiantes debe seguir en el centro de la educación, que las máquinas no pueden
desplazar a los docentes y que sus funciones y competencias deben revisarse para prepararlos
a trabajar en entornos con inteligencia artificial \citep[párrs.~12--13]{unesco2019beijing}.
Sobre la comparación, pide fomentar la investigación comparada entre países
\citep[párr.~31]{unesco2019beijing}. El Consenso no es vinculante: invita a los gobiernos a
considerar sus recomendaciones de acuerdo con su legislación y sus prácticas
\citep[p.~4]{unesco2019beijing}.

\subsubsection*{b) La Recomendación sobre la ética de la inteligencia artificial}

La Recomendación sobre la ética de la inteligencia artificial, adoptada por la Conferencia
General de UNESCO el 23 de noviembre de 2021, da el piso ético sobre el que se leen los
demás marcos. Se organiza en cuatro
valores, diez principios y once ámbitos de acción política \citep{unesco2021ai}. Los
valores son el respeto de los derechos humanos y la dignidad humana, la prosperidad del
medio ambiente y los ecosistemas, la diversidad y la inclusión, y la vida en sociedades
pacíficas, justas e interconectadas (párrs.~13--24). Los principios los hacen más
concretos (párr.~10) e incluyen la proporcionalidad, la equidad y no discriminación, la
protección de datos, la supervisión humana, la transparencia, la rendición de cuentas y,
de manera expresa, la sensibilización y la educación (párrs.~25--47).

La educación aparece dos veces. Es una de las esferas de competencia de UNESCO que la
Recomendación atiende en especial, porque vivir en sociedades digitalizadas exige nuevas
prácticas educativas, reflexión ética, pensamiento crítico y nuevas competencias
(párr.~3). Y es el octavo ámbito de acción política, ``Educación e investigación''
(párrs.~101--111). Ese ámbito pide alfabetización en inteligencia artificial para toda la
población (párr.~101); investigación sobre el uso responsable de la inteligencia artificial
en la enseñanza y en la formación de docentes, con la advertencia de que debe apoyar el
aprendizaje sin reducir las capacidades cognitivas y sin recabar información sensible
(párr.~104); y planes de estudio sobre ética de la inteligencia artificial para todos los
niveles, de acuerdo con las tradiciones educativas de cada país y en las lenguas locales
(párr.~106). Para esta tesis, la Recomendación funciona como referencia para el apartado 6.4
de la guía, el de equidad, inclusión y ética, y como contrapeso a su lectura
tecno-solucionista.

\subsubsection*{c) Los siete apartados de la guía de UNESCO como criterio de comparación}

El Consenso de Beijing encargó a UNESCO elaborar lineamientos para apoyar a los Estados
miembros en el diseño de políticas \citep[párr.~40]{unesco2019beijing}. La respuesta fue
\textit{AI and Education: Guidance for Policy-makers} \citep{miao2021guidance}, que se
presenta como una guía para desarrollar políticas de inteligencia artificial y educación y
ofrece recomendaciones para planear políticas en contextos locales
\citep[p.~1]{miao2021guidance}. Su sección 6 reúne esas recomendaciones en siete
apartados\footnote{Los títulos en español son traducción propia de la versión en inglés.}
\citep[pp.~31--37]{miao2021guidance}:

\begin{enumerate}[label=6.\arabic*]
  \item \textbf{Visión sistémica y prioridades estratégicas.} Definir una visión de la
  inteligencia artificial en todo el sistema educativo, evaluar la preparación del sistema
  y elegir prioridades. Propone cuatro metas: uso inclusivo y equitativo, mejora de la
  educación y el aprendizaje, desarrollo de competencias para la vida en la era de la
  inteligencia artificial, y uso transparente y auditable de los datos educativos.
  \item \textbf{Principio rector.} Adoptar un enfoque humanista, con la inteligencia
  artificial bajo control humano y al servicio de las personas, y mantener a las personas
  en el centro de la educación.
  \item \textbf{Planeación interdisciplinaria y gobernanza intersectorial.} Movilizar
  expertos de varias disciplinas y sectores, formar a quienes deciden y establecer
  mecanismos de coordinación con un enfoque de todo el gobierno.
  \item \textbf{Políticas y regulación para un uso equitativo, inclusivo y ético.} Fijar
  metas de inclusión, leyes de protección de datos que hagan los datos educativos visibles
  y auditables para docentes, estudiantes y familias, y revisar sesgos, incluido el de
  género.
  \item \textbf{Planes maestros para la gestión, la enseñanza, el aprendizaje y la
  evaluación.} Es el apartado más extenso. Incluye la gestión escolar, el aprendizaje
  centrado en el estudiante, el empoderamiento de los docentes, el aprendizaje a lo largo
  de la vida y los valores y competencias para la vida y el trabajo.
  \item \textbf{Pilotaje, monitoreo y evaluación, y construcción de evidencia.} Adoptar
  tecnología según prioridades educativas y no por novedad, verificar las afirmaciones de
  los proveedores y financiar investigación independiente.
  \item \textbf{Fomento de innovaciones locales.} Promover el desarrollo local de
  tecnologías de inteligencia artificial para la educación.
\end{enumerate}

Conviene precisar qué hace la guía y qué hago yo con ella. UNESCO no llama a estos
apartados ``siete áreas de política''. Los ordena como una visión (6.1), un principio
rector (6.2) y cinco grupos de recomendaciones (6.3 a 6.7) \citep[p.~31]{miao2021guidance}.
La decisión de tratarlos como siete categorías de comparación es mía, y la justifico por
tres razones. Primero, forman una lista cerrada y numerada, específica de educación, que
no construí a partir de mi propio corpus. Segundo, la guía se escribió para que los
gobiernos diseñaran sus políticas con ella, así que preguntar cuánto de cada apartado
aparece en un documento nacional es usarla en la dirección que sus autores previeron.
Tercero, la guía reconoce que la relación entre inteligencia artificial y educación se
desarrollará de maneras muy distintas según las circunstancias nacionales
\citep[p.~5]{miao2021guidance}, lo que abre la puerta a comparar esas diferencias sin
suponer que todos los países deban llegar al mismo lugar.

La guía también propone una distinción útil para leer los apartados 6.5 y 6.1: aprender
\textit{con} inteligencia artificial, aprender \textit{sobre} inteligencia artificial y
aprender para la colaboración entre humanos e inteligencia artificial
\citep[p.~13]{miao2021guidance}. Esta distinción se corresponde con la de
\citet{schiff2022education}: aprender con inteligencia artificial es inteligencia
artificial para la educación; aprender sobre ella y para colaborar con ella es educación
para la inteligencia artificial.

\subsection{La competencia docente en inteligencia artificial}
\label{subsec:mt-docente}

\subsubsection*{a) La alfabetización en inteligencia artificial}

\citet{long2020ailiteracy} definen la alfabetización en inteligencia artificial como un
conjunto de competencias que permite a las personas evaluar críticamente las tecnologías de
inteligencia artificial, comunicarse y colaborar con ellas de manera eficaz y usarlas como
herramienta en línea, en casa y en el trabajo. Llegaron a esa definición mediante una
revisión exploratoria de 150 documentos, y organizaron diecisiete competencias en torno a
cinco preguntas: qué es la inteligencia artificial, qué puede hacer, cómo funciona, cómo
debería usarse y cómo la percibe la gente. Ellos mismos presentan el marco como un punto de
partida para la conversación, no como un resumen exhaustivo del campo.

\citet{ng2021ailiteracy} revisaron treinta artículos publicados entre 2016 y 2021 y
organizaron la alfabetización en inteligencia artificial en cuatro aspectos inspirados en
la taxonomía de Bloom: conocer y comprender, usar y aplicar, evaluar y crear, y las
cuestiones éticas. No propusieron una definición propia; adoptaron la de Long y Magerko.
Su revisión mostró que casi todos los estudios se habían hecho con estudiantes de primaria
y secundaria, y que solo dos se ocupaban de docentes como población
\citep[p.~3]{ng2021ailiteracy}.

Para la educación básica, el grupo de trabajo conjunto de la Association for the
Advancement of Artificial Intelligence y la Computer Science Teachers Association propuso
cinco grandes ideas: las computadoras perciben el mundo mediante sensores; los agentes
mantienen modelos del mundo y los usan para razonar; las computadoras pueden aprender de
datos; lograr que los agentes interactúen con naturalidad con los humanos es un reto
considerable; y las aplicaciones de inteligencia artificial pueden tener efectos positivos
y negativos en la sociedad \citep[pp.~9797--9798]{touretzky2019k12ai}. El grupo no
elaboró un currículo propio, sino lineamientos que los currículos futuros deberían
cumplir \citep[p.~9798]{touretzky2019k12ai}.

Los tres trabajos comparten un rasgo que importa para esta tesis: definen qué debe saber
quien aprende, pero no qué necesita un docente para enseñarlo. El Consenso de Beijing
reconoce esa brecha cuando pide a UNESCO integrar las competencias en inteligencia
artificial a los marcos de competencias docentes \citep[párr.~42]{unesco2019beijing}.

\subsubsection*{b) El Marco de Competencias de IA para Docentes}

El Marco de competencias para docentes en materia de IA de UNESCO responde a esa brecha.
Se publicó en inglés en 2024 y en español en 2025, y se presenta como el primer marco
mundial de su tipo \citep[p.~6]{unesco2024aicft}. Parte de un diagnóstico: en 2022 solo
siete países habían desarrollado marcos o programas sobre inteligencia artificial para
docentes \citep[p.~13]{unesco2024aicft}.

El marco es una matriz de cinco ámbitos por tres niveles de progresión, que da quince
bloques de competencias \citep[p.~21]{unesco2024aicft}. Los ámbitos son una forma de
pensar centrada en el ser humano, la ética de la inteligencia artificial, los fundamentos y
las aplicaciones de la inteligencia artificial, la pedagogía de la inteligencia artificial y
la inteligencia artificial para el desarrollo profesional. Los niveles son adquirir,
profundizar y crear \citep[pp.~21--22]{unesco2024aicft}. El primer nivel equivale a una
alfabetización básica en inteligencia artificial para docentes \citep[pp.~24--25]{unesco2024aicft},
lo que conecta el marco con los trabajos de \citet{long2020ailiteracy} y
\citet{ng2021ailiteracy}. El marco advierte que el desarrollo de competencias depende del
contexto y no es jerárquico ni lineal: un docente puede estar en el nivel de profundizar en
un ámbito y en el de adquirir en otro \citep[pp.~21, 25]{unesco2024aicft}.

Tres de sus principios importan para la intervención. El primero protege el lugar del
docente: las herramientas de inteligencia artificial nunca deberían diseñarse para
reemplazar su responsabilidad legítima \citep[p.~18]{unesco2024aicft}. El segundo reparte
responsabilidades: los reguladores, los proveedores y las instituciones deben asumir la
gobernanza antes de exigir a los docentes que apliquen principios
\citep[p.~17]{unesco2024aicft}. El tercero reconoce que el marco no basta por sí solo,
porque su uso depende de la infraestructura, la regulación de datos, las políticas y el
desarrollo profesional \citep[p.~22]{unesco2024aicft}.

De este marco derivo el instrumento de autoevaluación que aplicaré antes y después de la
intervención con docentes. Lo tomo con la misma reserva que la guía de 2021: proviene de la
misma institución y comparte autor principal con la guía, y la crítica de
\citet{mochizuki2025ethics} a la figura del docente que necesita desarrollar competencias
se aplica también aquí. Por eso el instrumento medirá la competencia que el docente
reconoce en sí mismo, y no la tomará como prueba de que la adopción de la inteligencia
artificial sea deseable.

\subsection{Revisión de estudios sobre formación docente en inteligencia artificial}
\label{subsec:mt-formacion}

La investigación sobre la formación docente en inteligencia artificial es todavía escasa, y
esa escasez es en sí misma un hallazgo. \citet{zawacki2019systematic} revisaron 146
artículos sobre aplicaciones de inteligencia artificial en la educación superior publicados
entre 2007 y 2018, seleccionados de 2{,}656 identificados inicialmente. Encontraron que la
mayoría provenía de las ciencias de la computación y de las disciplinas científicas y
tecnológicas, que casi no había reflexión crítica sobre los riesgos, y que la conexión con
perspectivas pedagógicas era débil. El subtítulo de su revisión resume el problema: dónde
están los educadores.

La revisión de \citet{ng2021ailiteracy} llegó a una conclusión parecida desde la
alfabetización en inteligencia artificial. De los treinta estudios que analizaron, la
mayoría se hizo con estudiantes de primaria y secundaria y solo dos tuvieron a docentes como
población \citep[p.~3]{ng2021ailiteracy}. Cuatro artículos discutieron cómo los programas
de formación podrían preparar a docentes sin conocimientos previos para incorporar la
alfabetización en inteligencia artificial al currículo \citep[p.~7]{ng2021ailiteracy}.

Cuando se preguntó a los docentes, los resultados apuntaron en la misma dirección.
\citet{chounta2022teachers} encuestaron a 140 docentes de educación básica y media de
Estonia, un país que encabezaba entre 27 países europeos un índice de preparación para el
aprendizaje digital a lo largo de la vida. Encontraron que los docentes tenían un
conocimiento limitado sobre la inteligencia artificial y sobre cómo podría apoyarlos en su
práctica, pero la percibían como una oportunidad para la educación. Identificaron además
necesidades ligadas al contexto: los docentes veían en la inteligencia artificial una
herramienta para acceder a contenido en varias lenguas y adaptarlo
\citep[p.~725]{chounta2022teachers}. Si en uno de los sistemas educativos mejor preparados
de Europa el conocimiento docente sobre inteligencia artificial era limitado, no hay razón
para suponer que en México sea mayor.

La política internacional reconoció esta brecha antes de que la investigación la llenara.
El Consenso de Beijing pidió revisar las funciones y competencias docentes y desarrollar
programas para prepararlos a trabajar en entornos con inteligencia artificial
\citep[párr.~13]{unesco2019beijing}. La guía de 2021 pidió definir las habilidades que los
docentes necesitan, actualizar sus marcos y programas de formación, y ofrecer capacitación
y apoyo continuo \citep[p.~35]{miao2021guidance}. Y el Marco de competencias de UNESCO
partió de constatar que en 2022 solo siete países contaban con marcos o programas de
inteligencia artificial para docentes \citep[p.~13]{unesco2024aicft}.

De esta revisión se desprenden tres conclusiones para la tesis. La primera es que la
formación docente en inteligencia artificial es un campo más prescrito que estudiado: los
organismos internacionales la recomiendan, pero la evidencia empírica sobre cómo formar a
los docentes es reducida y se concentra en países con sistemas educativos de alto
desempeño. La segunda es que los estudios disponibles miden sobre todo percepciones y
conocimientos, no el cambio en la competencia después de una intervención, lo que justifica
un diseño con medición previa y posterior. La tercera es que la investigación sobre
inteligencia artificial en educación se ha hecho con poca participación de los docentes y
con poco sustento pedagógico, de modo que una intervención diseñada desde la pedagogía y
con los docentes como protagonistas atiende un vacío documentado.
%% TODO: ampliar con estudios de intervención en formación docente (América Latina).
